
\subsection{Summary}
For the rotation field $\bm{\theta}(\reflenx)$, we observe that the Saint~Venant solution generates sectional rotations through terms of the form $(\FrameBasis \times \mathbf{v}) \times \bm{r}$. 
The rotation must therefore include contributions from torsion $(a_\vartheta + c_\vartheta)\FrameBasis$, bending $-x(\FrameBasis \times \IesanLoadM)$, and flexure $-\tfrac{1}{2}x^2(\FrameBasis \times \bm{b}_{\Section})$.

The distribution of the rigid body parameters between $\bm{u}$ and $\bm{\theta}$ requires careful consideration. 
The translational freedom $\bm{\delta}_u$ naturally enters $\bm{u}(\reflenx)$ as the displacement at $\reflenx=0$. 

Assembling these components yields:

\begin{equation}
\begin{aligned}
\bm{u}(\reflenx) &= \reflenx\,(a_\varepsilon + \StrainGauge_{\varepsilon})\FrameBasis - \tfrac{1}{2}x^2\IesanLoadM - \tfrac{1}{6}x^3\bm{b}_{\Section} - \tfrac{1}{2}\reflenx^2\,(\FrameBasis\otimes\CentroidLocation)\bm{b}_{\Section} \\
\bm{\theta}(\reflenx) &= \StrainGauge_\vartheta \FrameBasis - \FrameBasis \times \bm{\StrainGauge}_{\Section} + x(a_\vartheta + c_\vartheta)\FrameBasis - x(\FrameBasis \times \IesanLoadM) - \tfrac{1}{2}x^2(\FrameBasis \times \bm{b}_{\Section})
\end{aligned}
\label{eq:sv-kinematics-ii}
\end{equation}
The 1D deformations \(\bm{\gamma}=\bm{u}'+\FrameBasis\times\bm{\theta}\) and \(\bm{\kappa}=\bm{\theta}'\) follow from \eqref{eq:sv-kinematics-ii}:
\[
\begin{aligned}
\bm{\gamma}(\reflenx)&= (a_\varepsilon - \reflenx\,\CentroidLocation\cdot\bm{b}_{\Section})\,\FrameBasis +  \bm{\StrainGauge}_{\gamma},\\[2pt]
\bm{\kappa}(\reflenx)&= 
(a_\vartheta+c_\vartheta)\,\FrameBasis
-\,(\FrameBasis\times\IesanLoadM)\;-\;x\,(\FrameBasis\times\bm{b}_{\Section}).
\end{aligned}
\]


\begin{Quote}
\begin{equation}
\begin{aligned}
\bm{u}(\reflenx)&= 
  \reflenx \,a_\varepsilon\,\FrameBasis 
 -\tfrac12 \reflenx^2\,\IesanLoadM
 \;-\;\tfrac16 \reflenx^3\,\bm{b}_{\Section}
 \;-\;\tfrac12 x^2\,(\FrameBasis\otimes\CentroidLocation)\bm{b}_{\Section},\\[2pt]
\bm{\theta}(\reflenx)&= 
  \reflenx\,(a_\vartheta+c_\vartheta)\,\FrameBasis
  -\,\reflenx\,(\FrameBasis\times\IesanLoadM)\;-\;\tfrac{1}{2} \reflenx^2\,(\FrameBasis\times\bm{b}_{\Section})\;-\FrameBasis\times\bm{\StrainGauge}_{\gamma}.
\end{aligned}
% \label{eq:sv-kinematics-ii}
\end{equation}
\end{Quote}

\begin{Quote}
\cite{romano2012torsion} use the notation \(
\bm{d}(\reflenx) = -\FrameBasis\times\bm{\kappa} = \IesanLoadM + x\bm{b}_{\Section}
\)
\end{Quote}

\begin{remark}[Implementation]
The choice \eqref{eq:sv-kinematics-ii} has the following ideal properties:
\begin{itemize}
    \item Thus the transverse shear is constant (\(\bm{\gamma}_\perp=\BetaFromB\bm{b}_{\Section}\)) and the bending curvature is linear in \(\reflenx\).
    \item The constant offset \(-\FrameBasis\times\bm{\StrainGauge}_{\gamma}\) generates a \emph{constant} shear deformation \(\bm{\gamma}\). 
    \item \(\bm{u}\) is cubic in \(\reflenx\), and consequently, can be exactly represented by standard \(C^1\) Hermite beam finite elements.
    \item Under pure flexure with \(a_\varepsilon = \eta_\varepsilon = 0\) and centroidal coordinates: 
    \[
    \bm{\gamma} = \GaugeShearB + \GaugeShearA
    \]
\end{itemize}
\end{remark}

